\فصل{نتیجه‌گیری}

\قسمت{پیش‌بینی مسیر و تشکیل آرایش}

شبکهٔ پارامتریک در دامنهٔ آزمون، مسیر و رویداد پایان را بازتولید کرد.
با این شبکه، برنامه‌های رهاسازی برای چند خانوادهٔ هدف طراحی شدند.
مثلث، حروف و دو نمونهٔ اردک در بازپخش مستقل تأیید شدند.
آرایش استادیومی ۴۸ ذره‌ای در هر دو نوع رهاسازی پذیرفته شد.
در مجموعهٔ بررسی‌شده، بزرگ‌ترین آرایش پذیرفته‌شده ۷۲ ذره داشت.
آزمون صدذره‌ای اهمیت کنترل فاصله را در طول حرکت نشان داد.
اعتبار این نتایج به هندسه و شرایط مدل محدود است.

\قسمت{پیش‌بینی مستقیم و طراحی معکوس}

عملگر جریان، پاسخ به پروفایل‌های ورودی را در دامنهٔ آزموده‌شده تقریب زد.
خطای مؤلفه‌های میدان با مرجع اجزای محدود سنجیده شد.
برای انتقال مغناطیسی، مسیر پیش‌بینی‌شده با حل مستقل مقایسه شد.
دو آزمون منتشرشده نیز دقت حل‌های عصبی ویژهٔ حرکت ذره را نشان دادند.

در بازسازی جریان، خطای مسیر و خطای میدان جداگانه ارزیابی شدند.
تعداد و پوشش مسیرها بر کیفیت بازسازی ورودی اثر گذاشتند.
در طراحی معکوس مغناطیسی، تحریک‌های مجاز با بازپخش مستقل ارزیابی شدند.
آزمون تقارن علامت، وجود فرمان‌های هم‌ارز برای یک مسیر را نشان داد.
